\par\Needspace{6\baselineskip}\medskip
\begin{center}\hypertarget{native-192-article143}{}\pdfbookmark[2]{143.}{article-143}143.\end{center}
Simulac angulus $z$ erutus est, statim habetur $r'$ per aequationem
\[
r'=\frac{R'\sin\delta'}{\sin z}\mathpunct{\cdot}
\]
Porro ex aequationibus $P=\dfrac{n''}{n}$ atque III elicimus
\[
\frac{n'r'}{n}=\frac{(P+a)R'\sin\delta'}{b\sin(z-\sigma)}
\]
\[
\frac{n'r'}{n''}=\frac{1}{P}\cdot\frac{n'r'}{n}\mathpunct{\cdot}
\]

\SourceLeaf{193}Iam ut formulas, secundum quas situs punctorum $C,C''$ e situ puncti $C'$
determinandus est, tali modo tractemus, ut ipsarum veritas generalis pro iis
quoque casibus, quos figura 4 non monstrat, statim eluceat, observamus, sinum
distantiae puncti $C'$ a circulo maximo $CB$ (positive sumtae in regione
superiori, negative in inferiori) aequalem fieri producto ex $\sin\epsilon''$
in sinum distantiae puncti $C'$ a $D''$ secundum directionem progressivam
mensuratae, adeoque $=-\sin\epsilon''\sin C'D''=-\sin\epsilon''\sin(z+A'D''-\delta')$;
perinde fit sinus distantiae puncti $C''$ ab eodem circulo maximo
$=-\sin\epsilon'\sin C''D'$. Manifesto autem iidem sinus sunt ut $\sin CC'$ ad
$\sin CC''$, sive ut $\dfrac{n''}{rr'}$ ad $\dfrac{n'}{rr''}\mathpunct{\raisebox{.45ex}{,}}$ sive ut $n''r''$ ad
$n'r'$. Statuendo itaque $C''D'=\zeta''$, habemus
\[
\tag*{V.}r''\sin\zeta''=\frac{n'r'}{n''}\cdot\frac{\sin\epsilon''}{\sin\epsilon'}\sin(z+A'D''-\delta').
\]
Prorsus simili modo statuendo $CD'=\zeta$ eruitur
\[
\tag*{VI.}r\sin\zeta=\frac{n'r'}{n}\cdot\frac{\sin\epsilon}{\sin\epsilon'}\sin(z+A'D-\delta')
\]
\[
\tag*{VII.}r\sin(\zeta+AD''-AD')=r''P\cdot\frac{\sin\epsilon}{\sin\epsilon''}\sin(\zeta''+A''D-A''D').
\]
Combinando aequationes V et VI cum sequentibus ex art. 139 transscriptis
\[
\tag*{VIII.}r''\sin(\zeta''-A''D'+\delta'')=R''\sin\delta''
\]
\[
\tag*{IX.}r\sin(\zeta-AD'+\delta)=R\sin\delta,
\]
quantitates $\zeta,\zeta'',r,r''$ ad normam art. 78 inde derivabuntur.
Qui calculus quo commodius absolvatur, formulas ipsas huc attulisse haud
ingratum erit. Statuatur
\[
\tag*{17)}\frac{R\sin\delta}{\sin(AD'-\delta)}=\varkappa
\]
\[
\tag*{18)}\frac{R''\sin\delta''}{\sin(A''D'-\delta'')}=\varkappa''
\]
\[
\tag*{19)}\frac{\cos(AD'-\delta)}{R\sin\delta}=\lambda
\]
\[
\tag*{20)}\frac{\cos(A''D'-\delta'')}{R''\sin\delta''}=\lambda''.
\]
Computus harum quantitatum, aut potius logarithmorum earum, a $P$ et $Q$
etiamnum independens, tamquam negotium \emph{quintum} et ultimum in
operationibus quasi praeliminaribus spectandum est, commodeque statim cum
computo ipsarum $a,b$ sive cum negotio quarto absolvitur, ubi fit
$a=\dfrac{\varkappa}{\varkappa''}\mathpunct{\cdot}$ — Faciendo \SourceLeaf{194}dein
\[
\begin{aligned}
\frac{n'r'}{n}\cdot\frac{\sin\epsilon}{\sin\epsilon'}\cdot\sin(z+A'D-\delta')&=p\\[.5ex]
\frac{n'r'}{n''}\cdot\frac{\sin\epsilon''}{\sin\epsilon'}\cdot\sin(z+A'D''-\delta')&=p''\\[.5ex]
\varkappa(\lambda p-1)&=q\\[.5ex]
\varkappa''(\lambda''p''-1)&=q''
\end{aligned}
\]
eliciemus $\zeta$ et $r$ ex $r\sin\zeta=p$, $r\cos\zeta=q$, atque $\zeta''$ et $r''$ ex $r''\sin\zeta''=p''$,
$r''\cos\zeta''=q''$. Ambiguitas in determinandis $\zeta$ et $\zeta''$ hic adesse nequit, quia $r$
et $r''$ necessario evadere debent quantitates positivae. Calculus perfectus per
aequationem VII si lubet confirmari poterit.

Sunt tamen duo casus, ubi aliam methodum sequi oportet. Quoties scilicet punctum $D'$ cum $B$ vel coincidit vel ipsi in sphaera oppositum est,
sive quoties $AD'-\delta=0$ vel $=180^0$, aequationes VI et IX necessario identicae esse debent, fieretque $\varkappa=\infty$, $\lambda p-1=0$, adeoque $q$ indeterminata. In
hoc casu $\zeta''$ et $r''$ quidem eo quo docuimus modo determinabuntur, dein vero
$\zeta$ et $r$ e combinatione aequationis VII cum VI vel IX elicere oportebit. Formulas ipsas ex art. 78 desumendas huc transscribere supersedemus; observamus tantummodo, quod in eo quoque casu, ubi est $AD'-\delta$ non quidem $=0$
neque $=180^0$, attamen arcus valde parvus, eandem methodum sequi praestat,
quoniam tunc methodus prior praecisionem necessariam non admitteret. Et
quidem adoptabitur combinatio aequationis VII cum VI vel cum IX, prout
$\sin(AD''-AD')$ maior vel minor est quam $\sin(AD''-\delta)$.

Perinde in casu, ubi punctum $D'$, vel ipsi oppositum, cum $B''$ vel coincidit vel parum ab eodem distat, determinatio ipsarum $\zeta'',r''$ per methodum
praecedentem vel impossibilis vel parum tuta foret. Tunc itaque $\zeta$ et $r$ quidem per illam methodum determinabuntur, dein vero $\zeta''$ et $r''$ e combinatione
aequationis VII vel cum V vel cum VIII, prout $\sin(A''D-A''D')$ maior vel
minor est quam $\sin(A''D-\delta'')$. Ceterum haud metuendum est, ne \emph{simul} $D'$
cum punctis $B,B''$ vel cum punctis oppositis coincidat, vel parum ab ipsis
distet: casum enim eum, ubi $B$ cum $B''$ coincidit, vel perparum ab eo distat,
iam supra art. 138 a disquisitione nostra exclusimus.
